Existence of the isotropic pair: in the orthonormal basis of Section \ref{sec:setting} take
\begin{align*}
w_1 &:= a_1 + b_1, & w_2 &:= a_2 + b_2, & w_3 &:= b_3 + I_7\,b_4,\\
w_1' &:= \tfrac12(a_1 - b_1), & w_2' &:= \tfrac12(a_2 - b_2), & w_3' &:= -\tfrac12(b_3 - I_7\,b_4).
\end{align*}
The first two isotropic vectors pair a plus direction with a minus direction and are real; the third pairs two minus directions and is complex, with $I_7 = I_6e$ the central unit of $\Cl(3,4) = \Cl(2,4)_\C$ (\claimref{std-algebra.complexification-is-the-odd-extension}): $Q(b_3 + I_7b_4) = Q(b_3) - Q(b_4) + 2\ip{b_3}{b_4}I_7 = -1 + 1 = 0$; and $\ip{w_j}{w_k'} = \delta_{jk}$, $\ip{w_j}{w_k} = \ip{w_j'}{w_k'} = 0$ by inspection. This is the basis the verification script uses, over the Gaussian rationals $\mathrm{span}_{\mathbb{Q}}\{1, I_7\}$. Conjugacy: two dual pairs $(W, W')$ and $(U, U')$ with dual bases $w_j, w_j'$ and $u_j, u_j'$ have the same Gram matrix $\bigl(\begin{smallmatrix} 0 & I \\ I & 0\end{smallmatrix}\bigr)$, so the linear map $w_j \mapsto u_j$, $w_j' \mapsto u_j'$ is an isometry of $V_\C$ carrying one pair to the other. The Clifford relation: for $w_1, w_2 \in W$, $\{w_1\wedge, w_2\wedge\} = 0$ and $\ip{w_1}{w_2} = 0$; for $w'_1, w'_2 \in W'$ likewise with $\{\iota_{w'_1}, \iota_{w'_2}\} = 0$; and the mixed case is the graded Leibniz identity: contraction is the degree $-1$ derivation extending $\iota_{w'}(w) = \ip{w}{w'}$, so for $\xi \in \Lambda^k W$,
\[
\iota_{w'}(w\wedge\xi) = \ip{w}{w'}\,\xi - w\wedge\iota_{w'}\xi, \qquad\text{that is}\qquad \{w\wedge, \iota_{w'}\} = \ip{w}{w'}\,\id ,
\]
which the script also verifies on every basis state. Expanding $u = w + w'$,
\[
c(u)^2 = 2\,(w\wedge + \iota_{w'})^2 = 2\,\{w\wedge, \iota_{w'}\} = 2\ip{w}{w'}\,\id = Q(u)\,\id ,
\]
using $Q(w+w') = 2\ip{w}{w'}$. Since $\Cl(V_\C) = T(V_\C)/\langle u\otimes u - Q(u)\rangle$ (\claimref{std-algebra.complex-clifford-algebra}(i), the complex form of \claimref{std-algebra.existence-uniqueness}), a linear map $c: V_\C \to A$ into an algebra with $c(u)^2 = Q(u)1$ factors uniquely through $\Cl(V_\C)$; so $c$ extends. Isomorphism: $\Cl(V_\C)$ is spanned by the $64$ ordered products of the six basis vectors (the Clifford relations reorder any word into one of them; \claimref{std-algebra.complex-clifford-algebra}(i)), and the script checks that their $64$ images are linearly independent in $\End(\Lambda^\bullet W)$; so $c$ is injective on a spanning set of size $64 = \dim \End(\Lambda^\bullet W)$, hence bijective, and $\Cl(V_\C) \cong \End(\C^8)$ is simple. Irreducibility and uniqueness: $\End(\C^8)$ acts irreducibly on $\C^8$ (any nonzero vector is carried onto any other by a rank-one map), and every module $M$ of a full matrix algebra is a direct sum of copies of $\C^8$: the matrix units give $M = \bigoplus_i E_{ii}M$ with $E_{1i}: E_{ii}M \to E_{11}M$ and $E_{i1}$ inverse isomorphisms, so $M \cong \C^8 \otimes E_{11}M$, and an irreducible $M$ is $\C^8$. Vectors are odd operators (wedge and contraction each shift exterior degree by one), so the even subalgebra preserves the parity split, and on an 8-dimensional module with a 4+4 split the even part $\Cl^0(V_\C)$, of dimension 32, maps onto $\End(\CS)\oplus\End(\CSm)$, of dimension $16 + 16 = 32$, injectively by simplicity of its two factors and the disjointness of their supports; the script checks the block structure exactly.
